\chapter{Time series and stationarity}\label{ch:ts}
A price series is not a dataset of independent draws --- it is an
ordered path. That single fact is why the tests in this book need
care: most of classical statistics assumes the thing you feed it is
\emph{stationary}, and prices emphatically are not. This chapter
makes that precise: what stationarity means, the two ways a series
fails it, the differencing fix you already used in ch1, and the ADF
test that tells you whether the fix was needed --- including the one
thing ADF cannot tell you.

\section{Stationarity, weak and strong}\label{sec:stat}
\index{stationarity}\index{weak stationarity}
A process is \emph{strongly} (strictly) stationary if the joint
distribution of any finite window $(x_{t_1},\dots,x_{t_k})$ equals
that of $(x_{t_1+h},\dots,x_{t_k+h})$ for every shift $h$ ---
statistics are translation-invariant in time, full stop. That is
stronger than anything you can check on one path, so practice uses
the weaker version: \emph{weakly} (covariance) stationary means
\begin{itemize}
\item $\mathbb E[x_t]=\mu$ constant,
\item $\operatorname{Var}[x_t]=\sigma^2$ constant, finite,
\item $\operatorname{Cov}[x_t,x_{t-k}]$ depends only on the lag $k$,
not on $t$.
\end{itemize}
Why tests need it: a sample mean $\bar x$ over a window is an
estimate of \emph{something} only if the mean doesn't move while you
average. On a random walk, $\mathbb E[x_t]$ is still defined but
$\operatorname{Var}[x_t]=t\sigma^2$ grows --- the second condition
fails, and with it every ``standard error'' computed under the
i.i.d.\ assumption.

\begin{example}[Identify by eye, then by number]\label{ex:tsgrid}
Four generators, same $T=500$:
white noise, AR(1) with $\phi=0.85$, Ornstein--Uhlenbeck with
$\theta=0.12$ around $\mu=4$, and a random walk. Which are
stationary?

\textbf{Strategy.} Eyeball for a constant level and a constant
spread, then confirm with the printed summaries.

\textbf{Solution.} Figure~\ref{fig:tsgrid} shows all four; real
output:
\begin{lstlisting}[numbers=none]
white noise : mean=+0.055 sd=1.001 first50 sd=1.092 last50 sd=0.978
AR(1) phi.85: mean=+0.305 sd=1.951
random walk : mean=+7.883 sd=9.361 start=+2.04 end=+27.50
OU theta.12 : mean=+4.011 sd=0.729 -> half-life 5.8 bars
\end{lstlisting}
White noise and OU look flat-and-constant --- stationary. AR(1) with
$\phi=0.85$ wanders more but is still stationary (an AR(1) with
$|\phi|<1$ has a finite stationary variance
$\sigma_\varepsilon^2/(1-\phi^2)$; the wandering is just slow
reversion, not a unit root). The random walk is the odd one out:
level drifts from $+2$ to $+27.5$ and the spread grows --- it fails
both the mean and variance conditions.
\end{example}

\begin{figure}[ht]\centering
\begin{tikzpicture}
\begin{axis}[width=0.48\linewidth,height=4.2cm,name=a,
  title={white noise},grid=major,grid style={gray!25},
  ymin=-3.2,ymax=3.2,title style={font=\small},
  xlabel style={font=\small},ylabel style={font=\small}]
\addplot[ink] coordinates {(0,2.041) (5,-0.216) (10,0.226) (15,-0.391) (20,0.024) (25,0.541) (30,-0.886) (35,0.670) (40,0.276) (45,-0.053) (50,-0.206) (55,-0.783) (60,-1.639) (65,-0.914) (70,0.328) (75,-1.048) (80,1.682) (85,-0.639) (90,-0.596) (95,-1.730) (100,-0.317) (105,0.134) (110,1.499) (115,0.533) (120,-0.979) (125,-0.197) (130,1.369) (135,0.658) (140,-0.931) (145,-1.171) (150,0.333) (155,0.111) (160,0.853) (165,1.695) (170,0.838) (175,-0.283) (180,1.019) (185,0.735) (190,-0.663) (195,-1.190) (200,1.632) (205,0.760) (210,-0.068) (215,0.042) (220,-2.138) (225,-1.951) (230,-0.462) (235,-1.231) (240,-0.080) (245,-0.788) (250,0.514) (255,0.160) (260,0.034) (265,0.753) (270,0.594) (275,-0.356) (280,-0.136) (285,-0.813) (290,-1.534) (295,-0.082) (300,0.190) (305,-1.696) (310,1.527) (315,0.881) (320,-0.471) (325,-1.108) (330,0.926) (335,1.103) (340,1.486) (345,0.203) (350,0.839) (355,-0.337) (360,0.133) (365,-0.890) (370,0.580) (375,0.754) (380,-0.519) (385,0.581) (390,0.299) (395,-0.380) (400,-1.332) (405,1.429) (410,0.543) (415,2.024) (420,0.452) (425,1.019) (430,1.038) (435,-0.587) (440,-1.774) (445,-0.406) (450,-0.029) (455,1.169) (460,0.853) (465,-0.083) (470,0.014) (475,0.786) (480,0.857) (485,-0.030) (490,0.827) (495,1.077) (499,0.859)};
\addplot[rulegray,dashed,domain=0:500]{0};
\end{axis}
\begin{axis}[width=0.48\linewidth,height=4.2cm,
  at={(a.right of east)},anchor=left of west,xshift=0.5cm,name=b,
  title={AR(1) $\phi=0.85$},grid=major,grid style={gray!25},
  ymin=-5.5,ymax=5.5,title style={font=\small}]
\addplot[accent] coordinates {(0,0.000) (5,-2.088) (10,0.302) (15,-1.993) (20,-0.233) (25,1.058) (30,1.104) (35,2.047) (40,-1.776) (45,-1.504) (50,1.405) (55,-0.046) (60,-2.155) (65,-3.613) (70,-2.177) (75,-1.915) (80,-0.177) (85,1.257) (90,-1.199) (95,-2.938) (100,-0.147) (105,-0.012) (110,4.099) (115,1.117) (120,-1.287) (125,1.044) (130,4.230) (135,2.545) (140,-0.485) (145,-0.273) (150,0.460) (155,-1.031) (160,-1.181) (165,2.000) (170,0.926) (175,2.812) (180,4.123) (185,3.548) (190,2.199) (195,-1.361) (200,2.387) (205,0.514) (210,0.592) (215,-0.464) (220,-3.293) (225,-2.341) (230,-1.807) (235,-3.136) (240,-2.087) (245,-1.647) (250,1.143) (255,2.296) (260,1.829) (265,1.875) (270,3.227) (275,2.597) (280,2.328) (285,1.007) (290,-0.205) (295,1.576) (300,-0.671) (305,-3.709) (310,0.505) (315,-0.534) (320,-0.297) (325,-3.748) (330,1.574) (335,1.623) (340,-0.028) (345,-0.669) (350,1.398) (355,2.637) (360,1.236) (365,-2.371) (370,-1.631) (375,0.098) (380,0.454) (385,1.904) (390,0.678) (395,0.433) (400,-1.099) (405,1.466) (410,4.094) (415,5.046) (420,2.476) (425,2.352) (430,3.035) (435,-1.458) (440,-2.172) (445,-2.430) (450,-1.506) (455,-0.174) (460,1.926) (465,1.005) (470,3.247) (475,0.248) (480,-0.944) (485,-1.255) (490,-0.742) (495,3.605) (499,3.022)};
\addplot[rulegray,dashed,domain=0:500]{0};
\end{axis}
\begin{axis}[width=0.48\linewidth,height=4.2cm,
  at={(a.below south)},anchor=above north,yshift=-0.8cm,name=c,
  title={OU $\theta=0.12$ (mean $\mu=4$)},xlabel={bar},
  grid=major,grid style={gray!25},ymin=1,ymax=7,
  title style={font=\small}]
\addplot[good] coordinates {(0,4.000) (5,4.028) (10,3.081) (15,4.835) (20,4.842) (25,4.591) (30,4.308) (35,3.087) (40,3.570) (45,3.367) (50,3.198) (55,3.018) (60,3.353) (65,3.919) (70,3.902) (75,3.209) (80,3.827) (85,4.454) (90,4.077) (95,5.211) (100,5.299) (105,5.442) (110,4.846) (115,4.286) (120,4.272) (125,3.760) (130,4.034) (135,4.817) (140,4.253) (145,3.966) (150,4.151) (155,5.082) (160,4.320) (165,4.586) (170,4.357) (175,4.241) (180,2.906) (185,3.113) (190,2.926) (195,4.033) (200,4.167) (205,4.232) (210,4.439) (215,4.235) (220,4.025) (225,4.377) (230,3.692) (235,3.777) (240,3.040) (245,3.451) (250,3.933) (255,1.825) (260,3.250) (265,5.180) (270,5.363) (275,5.215) (280,3.792) (285,4.234) (290,3.446) (295,3.635) (300,4.168) (305,3.713) (310,3.587) (315,2.923) (320,4.624) (325,5.594) (330,5.324) (335,4.378) (340,3.467) (345,3.592) (350,3.837) (355,3.601) (360,3.755) (365,3.819) (370,3.723) (375,3.549) (380,4.045) (385,4.565) (390,5.264) (395,5.004) (400,4.171) (405,3.624) (410,3.964) (415,4.199) (420,4.099) (425,3.388) (430,4.660) (435,3.208) (440,2.494) (445,2.667) (450,2.406) (455,2.931) (460,4.234) (465,3.885) (470,3.394) (475,3.818) (480,4.249) (485,3.860) (490,4.198) (495,3.511) (499,4.589)};
\addplot[amber,dashed,domain=0:500]{4};
\end{axis}
\begin{axis}[width=0.48\linewidth,height=4.2cm,
  at={(b.below south)},anchor=above north,yshift=-0.8cm,name=d,
  title={random walk},xlabel={bar},
  grid=major,grid style={gray!25},ymin=-10,ymax=30,
  title style={font=\small}]
\addplot[amber] coordinates {(0,2.041) (5,-1.333) (10,-0.901) (15,-3.649) (20,-2.623) (25,-0.680) (30,0.858) (35,2.791) (40,-1.368) (45,-2.215) (50,1.038) (55,0.281) (60,-2.441) (65,-5.536) (70,-5.998) (75,-6.901) (80,-5.926) (85,-3.571) (90,-5.859) (95,-8.359) (100,-6.494) (105,-6.287) (110,-1.581) (115,-2.540) (120,-4.859) (125,-3.125) (130,1.669) (135,1.747) (140,-0.158) (145,0.814) (150,1.196) (155,-0.422) (160,-1.981) (165,0.612) (170,-0.131) (175,3.473) (180,6.760) (185,9.428) (190,10.156) (195,7.697) (200,11.296) (205,10.637) (210,10.578) (215,9.414) (220,5.930) (225,5.014) (230,4.056) (235,0.749) (240,0.215) (245,-0.198) (250,2.504) (255,5.117) (260,6.243) (265,6.882) (270,10.264) (275,12.054) (280,14.234) (285,14.364) (290,14.509) (295,16.640) (300,15.054) (305,10.577) (310,12.742) (315,11.198) (320,11.400) (325,6.620) (330,10.066) (335,10.984) (340,8.813) (345,7.659) (350,10.084) (355,12.461) (360,12.306) (365,9.283) (370,8.799) (375,9.471) (380,10.342) (385,12.346) (390,12.833) (395,13.489) (400,12.243) (405,14.699) (410,19.596) (415,23.346) (420,23.100) (425,24.763) (430,27.455) (435,23.502) (440,22.027) (445,19.955) (450,19.097) (455,19.610) (460,22.550) (465,22.929) (470,27.045) (475,24.681) (480,22.647) (485,21.480) (490,20.550) (495,25.834) (499,27.497)};
\addplot[rulegray,dashed,domain=0:500]{0};
\end{axis}
\end{tikzpicture}

\caption{Which are stationary? White noise and OU are (constant mean
and spread); AR(1) with $\phi<1$ is stationary but lazy about
reverting; the random walk drifts and fans out.}
\label{fig:tsgrid}
\end{figure}

\begin{checkpoint}[Name the failure]
For each, say which weak-stationarity condition fails: (a)
$x_t=0.5t+\varepsilon_t$; (b) $x_t=x_{t-1}+\varepsilon_t$; (c)
$\varepsilon_t\cdot t$ with $\varepsilon$ i.i.d. ((a) mean, (b)
variance, (c) variance.)
\end{checkpoint}

\section{ACF, PACF and the Ljung--Box check}\label{sec:acfpacf}
\index{PACF}\index{Ljung--Box}
Ch1 used the ACF as a fact-detector. For model identification the
companion is the \emph{PACF}: $\operatorname{PACF}(k)$ is the
correlation of $x_t$ with $x_{t-k}$ \emph{after removing} the effect
of the intervening lags. The signature table worth memorising:
\begin{itemize}
\item White noise: both die immediately.
\item AR($p$): ACF decays smoothly, PACF cuts off after $p$.
\item MA($q$): ACF cuts off after $q$, PACF decays smoothly.
\item Unit root: ACF hugs 1 for dozens of lags
(Figure~\ref{fig:acf3}).
\end{itemize}
ACF bars are pointwise; the joint question ``is \emph{any} of the
first $m$ lags nonzero?'' is the Ljung--Box test:
$Q=m\,(m+2)\sum_{k\le m}\hat\rho_k^2/(n-k)$, $\chi^2_m$ under the
null of whiteness.

\begin{example}[Ljung--Box before and after fitting]\label{ex:lb}
Run the check on the raw AR(1) series, then on its residuals after
fitting (\texttt{code/ljung\_box.py}).

\textbf{Strategy.} Raw series should fail whiteness badly; residuals
of a correct fit should pass it. That before/after pair is the
diagnostic loop in miniature.

\textbf{Solution.} Real output:
\begin{lstlisting}[numbers=none]
white noise    LB(10)=   5.10 p=0.8847 -> consistent with white noise
AR(1) .85 raw  LB(10)=1385.33 p=0.0000 -> autocorrelation remains
AR(1) resid    LB(10)=   5.05 p=0.8880 -> consistent with white noise
(phi_hat=0.861 vs planted 0.85)
\end{lstlisting}
The raw AR(1) fails at $Q=1385$; the fitted residuals pass at
$Q=5.05$ --- nothing left for the model to explain. White noise
itself passes, as it must. Read the pattern, not just the p-value:
LB is the ``model is finished'' certificate, and it is checked on
\emph{residuals}, never on the raw series.
\end{example}

\begin{figure}[ht]\centering
\begin{tikzpicture}
\begin{axis}[width=0.32\linewidth,height=4.6cm,name=p1,
  title={ACF: random walk},xlabel={lag},grid=major,
  grid style={gray!25},xmin=0.5,xmax=15.5,ymin=-0.1,ymax=1.05,
  title style={font=\small},xlabel style={font=\small}]
\addplot[ybar,fill=amber,draw=none,bar width=2.5pt] coordinates {(1,0.990) (2,0.979) (3,0.968) (4,0.957) (5,0.947) (6,0.936) (7,0.926) (8,0.916) (9,0.905) (10,0.896) (11,0.887) (12,0.878) (13,0.869) (14,0.859) (15,0.849)};
\addplot[rulegray,dashed,domain=0.5:15.5]{0.09};
\addplot[rulegray,dashed,domain=0.5:15.5]{-0.09};
\end{axis}
\begin{axis}[width=0.32\linewidth,height=4.6cm,
  at={(p1.right of east)},anchor=left of west,xshift=0.4cm,name=p2,
  title={ACF: AR(1) $\phi=0.85$},xlabel={lag},grid=major,
  grid style={gray!25},xmin=0.5,xmax=15.5,ymin=-0.1,ymax=1.05,
  title style={font=\small},xlabel style={font=\small}]
\addplot[ybar,fill=accent,draw=none,bar width=2.5pt] coordinates {(1,0.858) (2,0.744) (3,0.643) (4,0.560) (5,0.489) (6,0.417) (7,0.360) (8,0.300) (9,0.238) (10,0.187) (11,0.160) (12,0.121) (13,0.078) (14,0.040) (15,0.008)};
\addplot[rulegray,dashed,domain=0.5:15.5]{0.09};
\addplot[rulegray,dashed,domain=0.5:15.5]{-0.09};
\end{axis}
\begin{axis}[width=0.32\linewidth,height=4.6cm,
  at={(p2.right of east)},anchor=left of west,xshift=0.4cm,
  title={ACF: OU $\theta=0.12$},xlabel={lag},grid=major,
  grid style={gray!25},xmin=0.5,xmax=15.5,ymin=-0.1,ymax=1.05,
  title style={font=\small},xlabel style={font=\small}]
\addplot[ybar,fill=good,draw=none,bar width=2.5pt] coordinates {(1,0.872) (2,0.774) (3,0.674) (4,0.582) (5,0.500) (6,0.413) (7,0.334) (8,0.253) (9,0.174) (10,0.110) (11,0.060) (12,0.021) (13,-0.003) (14,-0.035) (15,-0.061)};
\addplot[rulegray,dashed,domain=0.5:15.5]{0.09};
\addplot[rulegray,dashed,domain=0.5:15.5]{-0.09};
\end{axis}
\end{tikzpicture}

\caption{ACF signatures. Random walk: decays too slowly to see in 15
lags. AR(1): clean geometric decay $\approx0.85^k$. OU: same decay
family (it is an AR(1) around a mean) and that is the point --- the
ACF cannot tell AR(1)-mean-reverting from AR(1)-persistent; level
and the ADF regression carry that information.}
\label{fig:acf3}
\end{figure}

\section{Deterministic trend vs unit root}\label{sec:trend}
\index{unit root}\index{trend-stationary}
``Non-stationary'' has two flavours and they are not
interchangeable:
\begin{itemize}
\item \textbf{Trend-stationary:} $x_t=\alpha+\beta
t+\text{stationary}$. Remove the deterministic line and it is
stationary. Shocks die out.
\item \textbf{Difference-stationary (unit root):}
$x_t=x_{t-1}+\varepsilon_t$ (a random walk) or with drift. Every
shock is permanent; $\operatorname{Var}$ grows with $t$. Differencing
restores stationarity.
\end{itemize}
The distinction matters because the fix differs: detrending a unit
root leaves a non-stationary remainder (the line you removed was
never the real level); differencing a trend-stationary series
\emph{over-differences} --- you manufactured an MA unit root of the
opposite sign. Most financial prices look like the second case.

\section{Differencing and logs: the price\texorpdfstring{$\to$}{->}returns fix}\label{sec:diff}
\index{differencing}
If $P_t$ contains a unit root, $r_t=\log P_t-\log P_{t-1}$ removes
it --- which is the operational reason ch1 modelled returns, not
prices. Verify on the real series.

\begin{example}[BTC: unit root in the level, none in the return]
\label{ex:btcadf}
On the cached Kraken daily closes (\texttt{data/btc\_daily.csv};
\texttt{code/btc\_stationarity.py}), run ADF on the log-price and on
its first difference.

\textbf{Strategy.} The question a trading researcher actually asks:
``is my model's input legal?'' Prices fail; returns pass.

\textbf{Solution.} Real output:
\begin{lstlisting}[numbers=none]
bars=721 log-price range [10.981,11.735]
log-price  ADF= -1.865 p=0.3487 crit5%=-2.866 -> unit root not rejected
log-return ADF=-27.621 p=0.0000 crit5%=-2.866 -> stationary
\end{lstlisting}
Figure~\ref{fig:diff} shows the before/after: the level wanders for
two years (and ADF agrees: $-1.865$ is far inside the $-2.866$
critical value, Figure~\ref{fig:adfsketch}); one difference kills it
and the return rejects at $-27.6$. This is why ``run the test on
prices'' is a trap question --- the test \emph{should} fail there.
\end{example}

\begin{figure}[ht]\centering
\begin{tikzpicture}
\begin{axis}[width=0.8\linewidth,height=4.6cm,name=lp,
  grid=major,grid style={gray!25},xmin=0,xmax=720,
  ymin=10.9,ymax=11.8,title={BTC log-price (Kraken daily)},
  title style={font=\small}]
\addplot[ink] coordinates {(0,11.125) (4,11.117) (8,11.183) (12,11.148) (16,11.209) (20,11.283) (24,11.400) (28,11.420) (32,11.501) (36,11.437) (40,11.480) (44,11.484) (48,11.512) (52,11.521) (56,11.569) (60,11.468) (64,11.482) (68,11.458) (72,11.449) (76,11.492) (80,11.440) (84,11.429) (88,11.548) (92,11.557) (96,11.558) (100,11.541) (104,11.498) (108,11.488) (112,11.484) (116,11.486) (120,11.461) (124,11.477) (128,11.355) (132,11.418) (136,11.410) (140,11.291) (144,11.335) (148,11.317) (152,11.339) (156,11.375) (160,11.324) (164,11.321) (168,11.258) (172,11.318) (176,11.351) (180,11.352) (184,11.445) (188,11.453) (192,11.473) (196,11.454) (200,11.543) (204,11.549) (208,11.543) (212,11.590) (216,11.586) (220,11.580) (224,11.560) (228,11.550) (232,11.603) (236,11.562) (240,11.558) (244,11.519) (248,11.585) (252,11.587) (256,11.593) (260,11.596) (264,11.675) (268,11.686) (272,11.678) (276,11.683) (280,11.681) (284,11.650) (288,11.642) (292,11.668) (296,11.708) (300,11.679) (304,11.635) (308,11.619) (312,11.605) (316,11.613) (320,11.614) (324,11.638) (328,11.658) (332,11.673) (336,11.637) (340,11.602) (344,11.641) (348,11.714) (352,11.717) (356,11.638) (360,11.605) (364,11.615) (368,11.617) (372,11.646) (376,11.608) (380,11.538) (384,11.548) (388,11.515) (392,11.442) (396,11.338) (400,11.378) (404,11.416) (408,11.440) (412,11.406) (416,11.418) (420,11.381) (424,11.380) (428,11.381) (432,11.379) (436,11.385) (440,11.421) (444,11.412) (448,11.422) (452,11.464) (456,11.408) (460,11.401) (464,11.401) (468,11.258) (472,11.114) (476,11.154) (480,11.129) (484,11.123) (488,11.132) (492,11.117) (496,11.104) (500,11.183) (504,11.132) (508,11.184) (512,11.218) (516,11.160) (520,11.173) (524,11.099) (528,11.110) (532,11.150) (536,11.191) (540,11.224) (544,11.243) (548,11.269) (552,11.266) (556,11.240) (560,11.287) (564,11.290) (568,11.299) (572,11.269) (576,11.255) (580,11.247) (584,11.201) (588,11.182) (592,11.021) (596,11.032) (600,11.068) (604,11.084) (608,11.065) (612,10.999) (616,10.999) (620,11.035) (624,11.058) (628,11.069) (632,11.082) (636,11.076) (640,11.085) (644,11.078) (648,11.058) (652,11.066) (656,11.082) (660,11.062) (664,11.051) (668,11.181) (672,11.272) (676,11.274) (680,11.260) (684,11.287) (688,11.273) (692,11.250) (696,11.245) (700,11.348) (704,11.340) (708,11.333) (712,11.347) (716,11.335) (720,11.330) (720,11.330)};
\end{axis}
\begin{axis}[width=0.8\linewidth,height=4.2cm,
  at={(lp.below south)},anchor=above north,yshift=-0.9cm,
  xlabel={bar},grid=major,grid style={gray!25},
  xmin=0,xmax=720,ymin=-0.075,ymax=0.075,
  title={BTC log-return},title style={font=\small}]
\addplot[accent] coordinates {(0,-0.009) (4,-0.005) (8,0.003) (12,-0.017) (16,0.025) (20,0.070) (24,0.002) (28,0.016) (32,-0.006) (36,0.027) (40,0.003) (44,0.034) (48,-0.022) (52,0.001) (56,0.009) (60,0.022) (64,0.016) (68,-0.006) (72,0.029) (76,0.030) (80,0.012) (84,0.045) (88,0.002) (92,0.002) (96,-0.004) (100,0.022) (104,-0.020) (108,0.001) (112,-0.005) (116,0.002) (120,0.011) (124,-0.005) (128,-0.006) (132,-0.011) (136,-0.026) (140,0.007) (144,0.006) (148,0.027) (152,0.011) (156,-0.000) (160,-0.002) (164,0.009) (168,0.012) (172,0.027) (176,-0.011) (180,-0.004) (184,-0.005) (188,0.003) (192,0.011) (196,0.004) (200,0.006) (204,-0.001) (208,0.017) (212,0.028) (216,0.019) (220,-0.020) (224,0.011) (228,0.015) (232,-0.000) (236,0.003) (240,0.001) (244,0.023) (248,-0.004) (252,-0.012) (256,-0.002) (260,0.016) (264,0.009) (268,0.001) (272,0.001) (276,-0.020) (280,-0.004) (284,-0.015) (288,0.007) (292,0.013) (296,-0.015) (300,-0.019) (304,0.019) (308,-0.008) (312,-0.010) (316,0.010) (320,0.004) (324,0.009) (328,-0.004) (332,-0.010) (336,-0.008) (340,0.001) (344,0.027) (348,0.014) (352,-0.009) (356,0.016) (360,-0.032) (364,-0.001) (368,0.005) (372,-0.023) (376,0.005) (380,-0.004) (384,0.021) (388,-0.040) (392,-0.018) (396,0.003) (400,0.015) (404,0.005) (408,-0.004) (412,0.011) (416,0.003) (420,-0.007) (424,0.007) (428,-0.005) (432,0.003) (436,-0.001) (440,0.021) (444,0.004) (448,0.024) (452,0.000) (456,-0.011) (460,-0.022) (464,-0.042) (468,-0.003) (472,-0.005) (476,-0.012) (480,0.023) (484,-0.011) (488,-0.009) (492,0.004) (496,0.021) (500,-0.036) (504,0.030) (508,-0.016) (512,-0.031) (516,-0.027) (520,-0.024) (524,0.013) (528,-0.001) (532,0.004) (536,0.008) (540,-0.007) (544,-0.017) (548,-0.006) (552,-0.006) (556,0.024) (560,0.018) (564,0.006) (568,-0.012) (572,-0.006) (576,0.003) (580,0.009) (584,0.003) (588,-0.050) (592,-0.008) (596,-0.002) (600,0.010) (604,-0.028) (608,0.009) (612,-0.004) (616,-0.018) (620,0.014) (624,-0.019) (628,-0.003) (632,-0.011) (636,0.004) (640,-0.012) (644,-0.020) (648,-0.011) (652,0.009) (656,0.002) (660,-0.003) (664,0.011) (668,0.069) (672,0.012) (676,-0.011) (680,-0.007) (684,0.001) (688,-0.017) (692,0.018) (696,0.042) (700,0.014) (704,-0.001) (708,0.006) (712,0.010) (716,-0.022) (719,0.005)};
\addplot[rulegray,dashed,domain=0:720]{0};
\end{axis}
\end{tikzpicture}

\caption{BTC log-price (top) vs its first difference, the log-return
(bottom). Non-stationary level in, stationary series out.}
\label{fig:diff}
\end{figure}

\begin{checkpoint}[Predict the verdict]
What do you expect ADF to report on (a) raw $P_t$, (b) $\log P_t$,
(c) $\Delta\log P_t$ for the same series? (Fail, fail, pass ---
logging compresses but does not remove the unit root.)
\end{checkpoint}

\section{The ADF test: what it says, what it cannot say}\label{sec:adf}
\index{ADF test}
The augmented Dickey--Fuller test fits the regression
\[
\Delta x_t = \alpha + \beta t + \gamma\, x_{t-1}
  + \textstyle\sum_{j}\delta_j\,\Delta x_{t-j} + e_t
\]
and tests $H_0:\gamma=0$ (a unit root) against $\gamma<0$
(stationary, mean-reverting). The statistic is the t-ratio on
$\gamma$, but the null distribution is \emph{not} the t-distribution
--- it has its own critical values ($\approx-2.87$ at 5\% with a
constant, $-3.44$ with a trend; read them off the output).

\begin{example}[ADF on the four canonical series]\label{ex:adf}
\texttt{code/adf\_demo.py} runs the test (statsmodels
\texttt{adfuller}, AIC lag selection) on everything built so far.

\textbf{Strategy.} Check the planted truth against the verdicts,
including the cointegrated pair.

\textbf{Solution.} Real output:
\begin{lstlisting}[numbers=none]
white noise  ADF=-22.971  p=0.0000  -> REJECT unit root (stationary)
AR(1) .85    ADF= -6.039  p=0.0000  -> REJECT unit root (stationary)
OU .12       ADF= -5.792  p=0.0000  -> REJECT unit root (stationary)
random walk  ADF= -0.349  p=0.9182  -> cannot reject (non-stationary)
coint sprd z ADF= -5.371  p=0.0000  -> REJECT unit root (stationary)
coint y=2x+z ADF= -1.029  p=0.7426  -> cannot reject (non-stationary)
\end{lstlisting}
Every planted verdict is recovered, including the pair: $y_t$
contains the walk's unit root, the spread $z_t$ does not.
\end{example}

\begin{example}[Where ADF goes blind: the power problem]
\label{ex:adfpow}
Same script, second block: AR(1) with $\phi=0.98$ at $T=500$, three
seeds:
\begin{lstlisting}[numbers=none]
seed  3: ADF=-2.644 p=0.084 not detected
seed  7: ADF=-2.496 p=0.116 not detected
seed 11: ADF=-3.979 p=0.002 stationary
\end{lstlisting}
$\phi=0.98$ \emph{is} stationary --- and two of three seeds fail to
reject. Near the boundary, the test has low power: it is a one-sided
statement. ``Cannot reject'' means the data are consistent with a
unit root, not that there is one. The practical rule: a borderline
result tells you to look at the ACF decay and the economics, not to
declare victory either way.
\end{example}

\begin{figure}[ht]\centering
\begin{tikzpicture}
\begin{axis}[width=0.7\linewidth,height=4.8cm,
  xlabel={ADF statistic},grid=major,grid style={gray!25},
  xmin=-5,xmax=3,ymin=0,ymax=0.5,hide y axis,
  xlabel style={font=\small}]
\addplot[ink,domain=-4.5:2.5,samples=120]
  {0.42*exp(-((x+1.4)^2)/(2*0.9^2))};
\addplot[fill=amber,opacity=0.30,draw=none,domain=-4.5:-2.867,samples=60]
  {0.42*exp(-((x+1.4)^2)/(2*0.9^2))} \closedcycle;
\draw[amber,dashed,thick] (axis cs:-2.867,0) -- (axis cs:-2.867,0.45)
  node[above,font=\small,text=amber] {$-2.867$ (5\% critical)};
\node[font=\small,text=slate,align=center] at (axis cs:-3.8,0.16)
  {reject\\$H_0$};
\draw[accent,ultra thick] (axis cs:-1.822,0.02) -- (axis cs:-1.822,0.34);
\node[font=\small,text=accent,anchor=west,align=left]
  at (axis cs:-1.7,0.3) {BTC log-price\\$-1.822$: inside, not rejected};
\end{axis}
\end{tikzpicture}

\caption{ADF rejection region (sketch of the null distribution with
a constant). The BTC log-price statistic $-1.865$ sits inside the
non-rejection zone; only the far left tail rejects the unit root.}
\label{fig:adfsketch}
\end{figure}

\begin{checkpoint}[Read the asymmetry]
An ADF run gives $p=0.02$. What do you conclude --- and what would
$p=0.20$ \emph{not} let you conclude? (Reject at 5\%: stationary by
this test. The large p proves nothing about a unit root; it says the
data are consistent with one.)
\end{checkpoint}

\section{Mean reversion vs random walk: where a testable edge could live}\label{sec:ou}
\index{Ornstein--Uhlenbeck}\index{half-life}\index{cointegration}
Stationarity earned by differencing is ``no structure left''. The
interesting stationary objects are the ones with structure: a spread
that keeps returning. The workhorse model is the
Ornstein--Uhlenbeck process
$\mathrm dx_t=\theta(\mu-x_t)\,\mathrm dt+\sigma\,\mathrm dW_t$ ---
an AR(1) pointed at a mean --- and its tradable version is the
\emph{cointegrated pair}: two individually non-stationary series
$y_t=\beta x_t+z_t$ whose linear combination $z_t$ is stationary.
Then the \emph{half-life} $t_{1/2}=\log 2/\theta$ is the expected
decay time of a spread excursion --- the number that sizes holding
periods.

\begin{example}[Half-life recovered from data]\label{ex:ouhl}
Plant OU spreads at three speeds, recover $\theta,\mu$ and
$t_{1/2}$ by the AR(1) regression $x_t=a+b\,x_{t-1}+e_t$
($\hat\theta=1-b$), \texttt{code/ou\_halflife.py}.

\textbf{Strategy.} The estimator is four lines; the question is how
well it reads back the truth at 4000 bars.

\textbf{Solution.} Real output:
\begin{lstlisting}[numbers=none]
planted theta=0.05 mu=2.00 hl= 13.9 | recovered: theta=0.044 hl= 15.6
planted theta=0.12 mu=2.00 hl=  5.8 | recovered: theta=0.110 hl=  6.3
planted theta=0.30 mu=2.00 hl=  2.3 | recovered: theta=0.290 hl=  2.4
\end{lstlisting}
Recovery is tight for fast reversion (2.3$\to$2.4) and looser for
slow (13.9$\to$15.6, $\sim$12\% high): slow reversion gives each
sample less independent information, exactly the ADF power problem
in another costume. Figure~\ref{fig:oup} shows the annotated path.
\end{example}

\begin{figure}[ht]\centering
\begin{tikzpicture}
\begin{axis}[width=0.8\linewidth,height=5.6cm,
  xlabel={bar},grid=major,grid style={gray!25},
  xmin=0,xmax=300,ymin=1.5,ymax=6,title style={font=\small}]
\addplot[good] coordinates {(0,4.000) (3,3.785) (6,4.175) (9,3.764) (12,3.706) (15,4.835) (18,4.642) (21,4.806) (24,4.825) (27,4.396) (30,4.308) (33,3.153) (36,3.204) (39,3.500) (42,3.927) (45,3.367) (48,3.532) (51,3.521) (54,3.024) (57,3.197) (60,3.353) (63,4.051) (66,3.616) (69,3.886) (72,3.510) (75,3.209) (78,3.075) (81,3.458) (84,3.735) (87,3.910) (90,4.077) (93,5.190) (96,4.658) (99,5.527) (102,5.018) (105,5.442) (108,4.982) (111,5.409) (114,5.269) (117,4.036) (120,4.272) (123,3.645) (126,3.486) (129,3.929) (132,4.138) (135,4.817) (138,4.371) (141,4.792) (144,3.775) (147,4.098) (150,4.151) (153,4.074) (156,5.197) (159,4.757) (162,4.439) (165,4.586) (168,4.262) (171,4.167) (174,4.107) (177,4.102) (180,2.906) (183,2.714) (186,3.093) (189,3.610) (192,3.629) (195,4.033) (198,4.223) (201,4.204) (204,4.103) (207,3.697) (210,4.439) (213,4.552) (216,4.144) (219,4.394) (222,4.718) (225,4.377) (228,3.557) (231,3.202) (234,3.465) (237,3.110) (240,3.040) (243,3.752) (246,3.345) (249,4.346) (252,3.382) (255,1.825) (258,2.817) (261,3.953) (264,4.821) (267,5.687) (270,5.363) (273,5.000) (276,5.129) (279,4.154) (282,3.966) (285,4.234) (288,3.710) (291,3.432) (294,3.005) (297,3.665) (299,3.981)};
\addplot[amber,dashed,domain=0:300]{4};
\node[anchor=west,font=\small,text=amber] at (axis cs:190,4.35) {$\mu=4$};
\draw[<->,accent,thick] (axis cs:250,3.0) -- (axis cs:255.8,3.0)
  node[midway,below,font=\small,text=accent] {$t_{1/2}\approx5.8$};
\end{axis}
\end{tikzpicture}

\caption{OU path ($\theta=0.12,\mu=4$): every excursion decays back
with planted half-life $\approx5.8$ bars.}
\label{fig:oup}
\end{figure}

\begin{example}[A synthetic cointegrated pair]\label{ex:coint}
$x_t$ a random walk, $y_t=2x_t+z_t$ with $z_t$ OU
($\theta=0.10$), \texttt{simulate\_ts.py} + \texttt{adf\_demo.py}.
Real output:
\begin{lstlisting}[numbers=none]
coint pair  : y=2x+z, spread z mean=+0.120 sd=0.731
              (x sd 8.93, y sd 18.02)
coint sprd z ADF= -5.371  p=0.0000  -> stationary
coint y=2x+z ADF= -1.029  p=0.7426  -> unit root not rejected
\end{lstlisting}
Both legs fail the ADF test individually; their combination $z_t$
passes (Figure~\ref{fig:coint}). That is the entire mean-reversion
thesis in one output block: the edge, if any, lives in the
stationary residual --- a number you can name, estimate and test ---
not in ``two prices that look correlated''.
\end{example}

\begin{figure}[ht]\centering
\begin{tikzpicture}
\begin{axis}[width=0.8\linewidth,height=4.6cm,name=lv,
  grid=major,grid style={gray!25},xmin=0,xmax=500,
  ymin=-55,ymax=15,title={$x_t$ (random walk) and $y_t=2x_t+z_t$},
  title style={font=\small},legend style={font=\small}]
\addplot[ink] coordinates {(0,0.287) (5,0.738) (10,0.185) (15,0.101) (20,3.527) (25,3.641) (30,5.165) (35,4.623) (40,4.256) (45,4.395) (50,2.914) (55,3.713) (60,1.660) (65,1.932) (70,1.110) (75,0.782) (80,1.964) (85,3.506) (90,3.766) (95,4.549) (100,2.397) (105,-0.106) (110,1.353) (115,1.263) (120,0.636) (125,1.223) (130,-0.597) (135,-3.671) (140,-5.184) (145,-3.817) (150,-6.310) (155,-6.517) (160,-5.666) (165,-5.892) (170,-6.444) (175,-4.963) (180,-6.580) (185,-6.444) (190,-7.918) (195,-7.299) (200,-7.149) (205,-8.450) (210,-11.866) (215,-14.499) (220,-16.388) (225,-17.842) (230,-19.896) (235,-17.710) (240,-18.537) (245,-19.195) (250,-21.382) (255,-19.363) (260,-20.293) (265,-19.128) (270,-19.713) (275,-20.082) (280,-20.766) (285,-23.667) (290,-24.149) (295,-25.204) (300,-23.400) (305,-22.640) (310,-21.095) (315,-19.827) (320,-21.046) (325,-20.805) (330,-17.562) (335,-16.685) (340,-15.327) (345,-12.870) (350,-15.284) (355,-13.365) (360,-9.476) (365,-7.055) (370,-8.126) (375,-8.969) (380,-6.115) (385,-5.237) (390,-4.427) (395,-5.037) (400,-6.981) (405,-6.942) (410,-5.211) (415,-0.760) (420,2.149) (425,4.932) (430,0.230) (435,-1.219) (440,-5.669) (445,-2.964) (450,-3.363) (455,-1.076) (460,0.081) (465,-1.136) (470,0.047) (475,-1.769) (480,-2.035) (485,-0.423) (490,1.542) (495,1.179) (499,1.304)};
\addlegendentry{$x_t$}
\addplot[accent] coordinates {(0,0.574) (5,0.713) (10,0.247) (15,1.352) (20,6.864) (25,6.663) (30,10.475) (35,9.602) (40,8.576) (45,10.019) (50,8.128) (55,8.671) (60,3.450) (65,3.428) (70,2.338) (75,2.386) (80,4.668) (85,7.134) (90,7.302) (95,9.714) (100,5.178) (105,0.208) (110,3.846) (115,3.837) (120,0.854) (125,1.543) (130,-2.266) (135,-8.381) (140,-10.353) (145,-7.183) (150,-12.114) (155,-12.839) (160,-10.881) (165,-12.253) (170,-12.367) (175,-9.721) (180,-13.479) (185,-13.447) (190,-16.330) (195,-15.471) (200,-14.896) (205,-17.354) (210,-23.638) (215,-29.346) (220,-32.835) (225,-35.111) (230,-39.430) (235,-34.992) (240,-36.435) (245,-38.202) (250,-42.716) (255,-39.289) (260,-41.132) (265,-38.756) (270,-39.489) (275,-39.570) (280,-41.734) (285,-47.626) (290,-48.515) (295,-50.889) (300,-47.977) (305,-47.149) (310,-43.069) (315,-40.284) (320,-42.144) (325,-41.112) (330,-34.100) (335,-32.495) (340,-30.058) (345,-24.883) (350,-28.867) (355,-25.533) (360,-18.220) (365,-13.527) (370,-16.059) (375,-17.929) (380,-12.643) (385,-10.489) (390,-7.463) (395,-9.729) (400,-13.058) (405,-13.352) (410,-10.479) (415,-1.706) (420,5.045) (425,11.040) (430,1.048) (435,-2.032) (440,-10.794) (445,-4.953) (450,-6.927) (455,-3.191) (460,-1.715) (465,-3.430) (470,-0.374) (475,-2.939) (480,-3.887) (485,-0.550) (490,3.376) (495,2.361) (499,2.572)};
\addlegendentry{$y_t$}
\end{axis}
\begin{axis}[width=0.8\linewidth,height=4.2cm,
  at={(lv.below south)},anchor=above north,yshift=-0.9cm,
  xlabel={bar},grid=major,grid style={gray!25},
  xmin=0,xmax=500,ymin=-2.6,ymax=2.6,
  title={spread $z_t=y_t-2x_t$},title style={font=\small}]
\addplot[good] coordinates {(0,0.000) (5,-0.763) (10,-0.122) (15,1.150) (20,-0.190) (25,-0.620) (30,0.144) (35,0.355) (40,0.063) (45,1.229) (50,2.300) (55,1.245) (60,0.130) (65,-0.436) (70,0.118) (75,0.823) (80,0.740) (85,0.123) (90,-0.229) (95,0.617) (100,0.383) (105,0.420) (110,1.140) (115,1.311) (120,-0.417) (125,-0.904) (130,-1.071) (135,-1.039) (140,0.014) (145,0.451) (150,0.507) (155,0.196) (160,0.451) (165,-0.469) (170,0.521) (175,0.204) (180,-0.320) (185,-0.559) (190,-0.493) (195,-0.873) (200,-0.598) (205,-0.454) (210,0.095) (215,-0.347) (220,-0.058) (225,0.573) (230,0.362) (235,0.427) (240,0.639) (245,0.189) (250,0.047) (255,-0.562) (260,-0.546) (265,-0.500) (270,-0.063) (275,0.595) (280,-0.203) (285,-0.292) (290,-0.217) (295,-0.480) (300,-1.177) (305,-1.869) (310,-0.879) (315,-0.630) (320,-0.052) (325,0.498) (330,1.024) (335,0.874) (340,0.596) (345,0.857) (350,1.701) (355,1.197) (360,0.732) (365,0.583) (370,0.193) (375,0.008) (380,-0.414) (385,-0.014) (390,1.392) (395,0.345) (400,0.904) (405,0.532) (410,-0.057) (415,-0.187) (420,0.746) (425,1.176) (430,0.588) (435,0.406) (440,0.543) (445,0.976) (450,-0.202) (455,-1.039) (460,-1.878) (465,-1.158) (470,-0.468) (475,0.600) (480,0.183) (485,0.297) (490,0.293) (495,0.003) (499,-0.037)};
\addplot[amber,dashed,domain=0:500]{0};
\end{axis}
\end{tikzpicture}

\caption{Cointegration. Both legs contain the unit root; the linear
combination $z_t=y_t-2x_t$ is stationary (ADF $-5.37$).}
\label{fig:coint}
\end{figure}

\begin{example}[Trend-stationary or difference-stationary?]\label{ex:trendvsdiff}
\textbf{Strategy.} Two candidate fixes compete for trending series:
subtract a fitted line, or difference. Run ADF on both candidates
and let the data vote.

\textbf{Solution.} \texttt{trend\_vs\_diff.py} detrends the log BTC
price linearly, runs ADF on the residuals, and compares with ADF on
the returns:
\begin{verbatim}
detrended log-price ADF stat -2.799  (crit5% -1.941)
differenced (returns) ADF stat -27.635  (crit5% -1.941)
trend slope -18.8% per yr, R2 of trend fit = 0.298
verdict: returns reject hard (-27.6); detrended only marginally
     (-2.80 vs -1.94) -- thin evidence for trend-stationarity
\end{verbatim}
Returns reject the unit root at $-27.6$ --- decisively. The
detrended residuals reject only \emph{marginally} ($-2.80$ vs
$-1.94$ critical): barely inside the 5\% band, and the verdict is
fragile to lag choice and window. Meanwhile the trend explains
only $R^2=0.30$: a real ink-on-the-chart slope that is statistically
thin, because a random walk's ``trend'' is an accumulation of luck,
not a deterministic schedule.

\textbf{Significance.} This is the concrete form of
\S\ref{sec:trend}'s distinction, on real data, with the honest
wrinkle: borderline ADF rejections on detrended levels are
\emph{weak} evidence of trend-stationarity, not proof. When the
detrended answer is marginal and the differenced answer is
decisive, difference. \emph{Estimated} drift and a unit root
coexist happily --- never let an $R^2$ of 0.30 talk you into a
deterministic trend.
\end{example}

\section*{Summary sheet}
\addcontentsline{toc}{section}{Summary sheet}
\begin{table}[ht]\centering
\begin{tabularx}{\linewidth}{@{}lYY@{}}
\toprule
\tablehead\thead{Object} & \thead{Definition} & \thead{When to use}\\
\midrule
Weak stationarity & const $\mu,\sigma^2$, lag-only cov &
the licence statistical tests need\\
ACF/PACF & raw vs partial autocorr & fingerprinting AR/MA/unit root\\
Ljung--Box $Q$ & joint test of first $m$ lags & ``are residuals
white?''\\
Unit root & $x_t=x_{t-1}+\varepsilon_t$ & the permanent-shock kind
of non-stationarity\\
Differencing & $\Delta x_t$ or $\Delta\log$ & kills a unit root,
not a deterministic trend\\
ADF test & $H_0$: unit root & says ``consistent with unit root'',
never proves it\\
OU/half-life & $\mathrm dx=\theta(\mu-x)\mathrm dt+\sigma\mathrm dW$,
$t_{1/2}=\log2/\theta$ & sizing mean-reversion holds\\
Cointegration & $y-\beta x$ stationary while legs are not &
where spread edges live\\
\bottomrule
\end{tabularx}
\caption{Chapter 2 toolkit.}
\end{table}

\section{Common mistakes}\index{common mistakes}
\begin{itemize}
\item \textbf{Testing prices for stationarity at all.} They will
fail; the question is whether the \emph{transformed} series passes.
Test the object you model.
\item \textbf{Reading ``cannot reject'' as ``there is a unit
root''.} ADF is one-sided: a big p is absence of evidence, not
evidence of absence --- Example~\ref{ex:adfpow} shows a stationary
series failing to reject.
\item \textbf{Ljung--Box on the raw series.} $Q$ is a residual
diagnostic; a raw AR(1) \emph{should} fail (LB$=1385$ in
Example~\ref{ex:lb}) and that failure is information, not an error.
\item \textbf{Detrending a unit root / differencing a trend.}
Removing a line from a walk leaves a walk; differencing a trend
manufactures an MA unit root. Diagnose which kind first
(\S\ref{sec:trend}).
\item \textbf{Believing the ACF distinguishes AR from OU.} Their
ACFs are the same family (Figure~\ref{fig:acf3}); the difference is
the level they revert to, not the decay shape.
\item \textbf{Trading ``cointegration'' on correlation alone.}
Correlated non-stationary legs can diverge forever; the ADF-passing
residual is the object that reverts.
\item \textbf{Detrending a unit root.} A fitted line through a
random walk produces residuals that still fail ADF
(Example~\ref{ex:trendvsdiff}) and \emph{fake} mean reversion ---
the residuals oscillate around the line because the walk wanders.
Difference first, model the level never.
\item \textbf{Reading ADF as ``the test''.} ADF has three regression
variants (no constant, constant, trend) and lag selection changes
the verdict. State both, or the number is not reproducible.
\end{itemize}

\section*{Exercises}\addcontentsline{toc}{section}{Exercises}
\begin{exercise}[easy]
Classify each as stationary/non-stationary and name the failing
condition: (a) $\varepsilon_t$ i.i.d.; (b) $x_t=3+x_{t-1}+\varepsilon_t$
(drift); (c) $x_t=\sin(2\pi t/24)+\varepsilon_t$; (d)
$x_t=0.5x_{t-1}+\varepsilon_t$.
\end{exercise}
\begin{solution}
(a) yes; (b) no (variance grows, mean moves); (c) no --- mean is
periodic, so $\mathbb E[x_t]$ is not constant; (d) yes
($|\phi|<1$).
\end{solution}
\begin{exercise}[easy]
Write the AR(1) stationary variance
$\sigma_\varepsilon^2/(1-\phi^2)$ and evaluate at $\phi=0.85$,
$\sigma_\varepsilon=1$. Explain the printed sd $1.95$.
\end{exercise}
\begin{solution}
$1/(1-0.7225)=3.60$, sd $1.90$ --- matches $1.951$ within sampling
error.
\end{solution}
\begin{exercise}[easy]
State the ADF null and the 5\% critical value (constant, no trend)
from the output block. Which side is rejection?
\end{exercise}
\begin{solution}
$H_0:\gamma=0$ (unit root); crit $\approx-2.87$; reject when the
statistic is \emph{more negative} (left tail).
\end{solution}
\begin{exercise}[easy]
Convert a planted half-life to $\theta$: $t_{1/2}=10$ bars.
\end{exercise}
\begin{solution}
$\theta=\log2/10\approx0.069$.
\end{solution}
\begin{exercise}[easy]
A series has ACF $_1=0.99$, ACF$_{10}=0.91$. Suggest the diagnosis
and the fix.
\end{exercise}
\begin{solution}
Near-unit-root; difference the series (or, for prices, use log
returns), then re-check.
\end{solution}
\begin{exercise}[medium]
Show $\operatorname{Var}[x_t]=t\sigma^2$ for the random walk
$x_t=\sum_{s\le t}\varepsilon_s$, and hence why $\bar x$ over a
window is not a mean estimator.
\end{exercise}
\begin{solution}
Sum of $t$ independent variances. The variance depends on $t$, so no
window-average is estimating a fixed $\mu$ --- weak stationarity
condition 2 fails.
\end{solution}
\begin{exercise}[medium]
OU autocorrelation: for the discrete form
$x_t=\mu+\phi(x_{t-1}-\mu)+e_t$ with $\phi=1-\theta$, show
$\operatorname{ACF}(k)=\phi^k$ and evaluate the lag where it first
drops below $0.1$ at $\theta=0.12$.
\end{exercise}
\begin{solution}
Iterating, $x_{t+k}-\mu=\phi^k(x_t-\mu)+\text{noise}$, so
corr$=\phi^k$. $\phi=0.88$: $0.88^k<0.1$ at $k=\log0.1/\log0.88
\approx18$.
\end{solution}
\begin{exercise}[medium]
Interpret the borderline ADF output: seed 3 in
Example~\ref{ex:adfpow} gives $p=0.084$ on a series you know is
AR(1) with $\phi=0.98$. State the two correct readings.
\end{exercise}
\begin{solution}
(1) At 5\% you cannot reject a unit root --- the data are consistent
with one. (2) That is not evidence the process is non-stationary;
power near $\phi=1$ is low, so treat the verdict as inconclusive and
lean on theory/ACF.
\end{solution}
\begin{exercise}[medium]
Build your own cointegrated pair with $\beta=3$, OU spread
$\theta=0.08$, on 800 bars; report whether ADF recovers stationarity
of the spread and non-stationarity of $y$.
\end{exercise}
\begin{solution}
Implementation exercise: reuse \texttt{coint\_pair} parameters;
expect spread rejects, $y$ does not (as in
Example~\ref{ex:coint}).
\end{solution}
\begin{exercise}[medium]
Ljung--Box: why does fitting the AR(1) first change the verdict from
$Q=1385$ to $Q=5$? What is the residual series estimating?
\end{exercise}
\begin{solution}
The model removed the autocorrelation it was designed to capture;
$e_t=x_t-\hat\phi x_{t-1}$ estimates the innovation, which by
construction is white. LB on residuals asks ``did anything survive
the model?''
\end{solution}
\begin{exercise}[hard]
Over-differencing: let $x_t=\alpha+\beta t+\varepsilon_t$
(trend-stationary). Compute $\Delta x_t$ and show it now contains an
MA(1) with unit coefficient --- an artificial non-invertibility.
\end{exercise}
\begin{solution}
$\Delta x_t=\beta+\varepsilon_t-\varepsilon_{t-1}$: a constant plus
an MA(1) at $\theta_1=-1$, i.e.\ on the unit circle in MA space ---
you created the pathology you were trying to remove.
\end{solution}
\begin{exercise}[hard]
Half-life estimation error: at fixed $T$, argue qualitatively why
$\hat\theta$ for small $\theta$ is biased, using the output ordering
in Example~\ref{ex:ouhl}. Connect to the ADF power problem.
\end{exercise}
\begin{solution}
Small $\theta$ means few effective reversions per sample (long
memory): the regression on $x_{t-1}$ is near-collinear, $b\uparrow1$,
$\hat\theta=1-b$ is squashed toward 0, so half-life overestimates ---
the same low-power regime as ADF near $\phi=1$.
\end{solution}
\begin{exercise}[hard]
Extend \texttt{btc\_stationarity.py}: split the 721-bar series in
half and run ADF on each half's returns. Report whether the verdict
is stable, and state the sampling caveat.
\end{exercise}
\begin{solution}
Implementation exercise; both halves should reject comfortably
(returns are far from the boundary), but treat each as a different
sample of a possibly time-varying process --- stationarity of
returns is a model statement, not a law.
\end{solution}
\begin{exercise}[hard]
A pairs trader reports ``correlation between the two log-prices is
0.97, so the spread is safe''. Refute using
Example~\ref{ex:coint}'s numbers.
\end{exercise}
\begin{solution}
Correlation of levels is meaningless under unit roots (both legs are
non-stationary and can drift apart without bound); the safety claim
requires the \emph{residual} $y-\beta x$ to pass ADF --- a different
and testable statement.
\end{solution}
\begin{exercise}[hard]
Implement a Hurst-style check on the cached BTC returns: compare
$\operatorname{sd}[\sum_{i\le k}r_i]$ with $k^{0.5}$ scaling on
$k\in\{1,4,16,64\}$. Do BTC daily returns scale like a random walk?
\end{exercise}
\begin{solution}
Implementation exercise; the fit is $\operatorname{sd}\propto k^H$
with $H\approx0.5$ for near-independent increments --- expect $H$
close to but not exactly $0.5$, and say what the deviation could
mean (mild serial structure vs heavy tails inflating the sd).
\end{solution}

\begin{exercise}[easy]
In Example~\ref{ex:trendvsdiff}, which ADF statistic is the
stronger evidence of stationarity, and why is the $-27.6$ not a
surprise?
\end{exercise}
\begin{solution}
$-27.6$ on returns --- returns of a unit-root series are white-ish
noise by construction, so decisive rejection is expected; the
detrended level rejects only marginally ($-2.80$ vs $-1.94$), which
is thin and fragile evidence for trend-stationarity.
\end{solution}
\begin{exercise}[medium]
Your ADF on a trending series gives stat $=-1.5$ (fail) with the
\texttt{"c"} variant but $-3.5$ (pass) with \texttt{"ct"}.
Write one sentence a reviewer could not attack.
\end{exercise}
\begin{solution}
``With a trend term the series is trend-stationary at 5\%, without
it it carries a unit root --- the two variants test different nulls,
so I report both and difference regardless''.
\end{solution}
\begin{exercise}[hard]
Show that if $y_t$ follows a pure random walk, the OLS $R^2$ of a
time-trend regression tends to a nondegenerate random variable, not
to $0$ --- i.e. trending R-squared values mean nothing. Simulate it
on 500-step walks to see the distribution.
\end{exercise}
\begin{solution}
A spurious-regression result (Granger--Newbold): the $t$-statistic
and $R^2$ both diverge/level off at random values because the
walk's low-frequency content mimics any trend. Simulation shows
$R^2$ spread over $[0.1, 0.9]$ with no point mass at 0.
\end{solution}

\section*{Notes}\addcontentsline{toc}{section}{Notes}
\textbf{Software pin.} \texttt{statsmodels 0.14.6} and
\texttt{scipy 1.13.1} are pinned in \texttt{code/README.md};
The \texttt{adfuller} defaults
are used verbatim (regression \texttt{"c"}, AIC lag choice). \textbf{Standing disclaimer.} This book teaches
research methods and market mechanics; nothing here is investment,
legal or tax advice. \textbf{Data.} The BTC series is the cached
Kraken daily file (\texttt{code/data/btc\_daily.csv}, fetched
2026-10-11); venue availability is \textbf{[unverified]} per
\texttt{docs/unverified.md}.
\textbf{Where this leads.} Ch3 drops to the mechanism that prints
those prices --- the order book --- and ch8 turns this chapter's
``is the input legal?'' into a harness rule.

\section*{Further reading}\addcontentsline{toc}{section}{Further reading}
\begin{itemize}
\item R.~Hyndman \& G.~Athanasopoulos, \emph{FPP3}, ch9
``Stationarity and differencing'' --- the pedagogy model for
\S\ref{sec:stat}--\S\ref{sec:diff}, free online:

\url{https://otexts.com/fpp3/stationarity.html}
\item MIT 18.S096 lecture notes 8 (Kempthorne) --- the course notes
behind the ARIMA/ADF treatment; free PDF:

\url{https://ocw.mit.edu/courses/18-s096-topics-in-mathematics-with-applications-in-finance-fall-2013/1926c83ecd7ea700f7cb63914c6d7c0f_MIT18_S096F13_lecnote8.pdf}
\item statsmodels \texttt{adfuller} documentation --- every option
and the critical-value table used in \S\ref{sec:adf}:

\url{https://www.statsmodels.org/stable/generated/statsmodels.tsa.stattools.adfuller.html}
\item ``Augmented Dickey--Fuller test'' --- compact reference for
the regression and the nonstandard null distribution:

\url{https://en.wikipedia.org/wiki/Augmented_Dickey\%E2\%80\%93Fuller_test}
\item MIT OCW 18.S096 course page --- context for the Watch list,
all lectures free:

\url{https://ocw.mit.edu/courses/18-s096-topics-in-mathematics-with-applications-in-finance-fall-2013/}
\item M.~L\'opez de Prado, lecture library --- the stationarity and
fractional-differentiation lectures that take this further:

\url{https://www.quantresearch.org/Lectures.htm}
\end{itemize}

\begin{rulebox}[Watch instead]
\begin{itemize}
\item \textbf{MIT OpenCourseWare --- ``8. Time Series Analysis I''}
(18.S096), 1:16:19 --- Kempthorne covers stationarity, Wold
representation and ARMA estimation; maps to
\S\ref{sec:stat}--\S\ref{sec:trend}. (Verified 2026-10-11.)
\url{https://www.youtube.com/watch?v=uBeM1FUk4Ps}
\item \textbf{MIT OpenCourseWare --- ``11. Time Series Analysis II''}
(18.S096), 1:23:47 --- multivariate series and least-squares
estimation; background for \S\ref{sec:ou}. (Verified 2026-10-11.)
\url{https://www.youtube.com/watch?v=_IFUfFuyQlU}
\end{itemize}
\end{rulebox}
